\section{Conclusion}

\label{sec:conclusion}
% ============================================================

We have presented a unified threshold cryptographic framework that covers 17 blockchain
families through two complementary protocols: FROST for Schnorr/EdDSA chains and
CGGMP21 for ECDSA chains. The framework provides:

\begin{enumerate}[nosep]
  \item \textbf{Universal coverage}: Every major blockchain is supported with native
    signature verification---no custom verifier contracts.
  \item \textbf{Constant on-chain cost}: One signature verification regardless of
    threshold size, yielding up to 90\% gas savings over multisig.
  \item \textbf{Strong security}: Unforgeability under $t{-}1$ corruption (Theorems
    \ref{thm:frost-uf-full},~\ref{thm:cggmp-uf}), identifiable abort
    (Theorem~\ref{thm:cggmp-ia}), forward security under proactive refresh
    (Theorem~\ref{thm:forward}), and safe protocol composition
    (Theorem~\ref{thm:composition}).
  \item \textbf{Bitcoin Taproot}: FROST KeygenTaproot produces x-only public keys
    (BIP-340) for key-path spends indistinguishable from single-signer outputs.
  \item \textbf{MPC group address}: On-chain verification uses a single aggregate
    key ($\texttt{ecrecover} == \texttt{mpcGroupAddress}$), not individual signers.
  \item \textbf{Operational safety}: Configurable rotation delay (default 24h, 0 for
    instant, max 7d) with on-chain visibility; proactive key refresh without public
    key change; LSS dynamic resharing for threshold changes.
  \item \textbf{Production performance}: 98ms end-to-end FROST signing, 47ms CGGMP21
    online signing with presignatures, 0.14ms LSS reshare ($3 \to 5$), in
    geographically distributed deployment.
\end{enumerate}

The framework is deployed in Lux Teleport, securing cross-chain asset transfers across
270+ blockchains. Implementation is open-source at \texttt{github.com/luxfi/threshold}
and \texttt{github.com/luxfi/mpc}.

% ============================================================
% References
% ============================================================
\bibliographystyle{plain}
\begin{thebibliography}{99}

\bibitem{komlo2020frost}
C.~Komlo and I.~Goldberg,
``FROST: Flexible Round-Optimized Schnorr Threshold Signatures,''
\textit{Selected Areas in Cryptography (SAC)}, 2020.
\url{https://eprint.iacr.org/2020/852}

\bibitem{canetti2021cggmp}
R.~Canetti, R.~Gennaro, S.~Goldfeder, N.~Makriyannis, and U.~Peled,
``UC Non-Interactive, Proactive, Threshold ECDSA with Identifiable Aborts,''
\textit{ACM CCS}, 2020. Revised 2021.
\url{https://eprint.iacr.org/2021/060}

\bibitem{bip340}
P.~Wuille, J.~Nick, and T.~Ruffing,
``BIP-340: Schnorr Signatures for secp256k1,''
Bitcoin Improvement Proposal, 2020.
\url{https://github.com/bitcoin/bips/blob/master/bip-0340.mediawiki}

\bibitem{bip341}
P.~Wuille, J.~Nick, and A.~Towns,
``BIP-341: Taproot: SegWit version 1 spending rules,''
Bitcoin Improvement Proposal, 2020.
\url{https://github.com/bitcoin/bips/blob/master/bip-0341.mediawiki}

\bibitem{bip342}
P.~Wuille, J.~Nick, and A.~Towns,
``BIP-342: Validation of Taproot Scripts,''
Bitcoin Improvement Proposal, 2020.
\url{https://github.com/bitcoin/bips/blob/master/bip-0342.mediawiki}

\bibitem{pedersen1991}
T.~P.~Pedersen,
``Non-Interactive and Information-Theoretic Secure Verifiable Secret Sharing,''
\textit{CRYPTO}, 1991.

\bibitem{feldman1987}
P.~Feldman,
``A Practical Scheme for Non-Interactive Verifiable Secret Sharing,''
\textit{FOCS}, 1987.

\bibitem{gennaro2018}
R.~Gennaro and S.~Goldfeder,
``Fast Multiparty Threshold ECDSA with Fast Trustless Setup,''
\textit{ACM CCS}, 2018.
\url{https://eprint.iacr.org/2019/114}

\bibitem{gennaro2020}
R.~Gennaro and S.~Goldfeder,
``One Round Threshold ECDSA with Identifiable Abort,''
IACR ePrint 2020/540, 2020.
\url{https://eprint.iacr.org/2020/540}

\bibitem{gennaro2003}
R.~Gennaro, S.~Jarecki, H.~Krawczyk, and T.~Rabin,
``Secure Distributed Key Generation for Discrete-Log Based Cryptosystems,''
\textit{Journal of Cryptology}, 2003.

\bibitem{nick2021musig2}
J.~Nick, T.~Ruffing, and Y.~Seurin,
``MuSig2: Simple Two-Round Schnorr Multi-Signatures,''
\textit{CRYPTO}, 2021.
\url{https://eprint.iacr.org/2020/1261}

\bibitem{ruffing2022roast}
T.~Ruffing, V.~Ronge, E.~C.~Fletcher, J.~Griffy, and D.~Schr\"oder,
``ROAST: Robust Asynchronous Schnorr Threshold Signatures,''
\textit{ACM CCS}, 2022.
\url{https://eprint.iacr.org/2022/550}

\bibitem{doerner2018}
J.~Doerner, Y.~Kondi, E.~Lee, and A.~Shelat,
``Secure Two-party Threshold ECDSA from ECDSA Assumptions,''
\textit{IEEE S\&P}, 2018.

\bibitem{pointcheval1996}
D.~Pointcheval and J.~Stern,
``Security Proofs for Signature Schemes,''
\textit{EUROCRYPT}, 1996.

\bibitem{bellare2006}
M.~Bellare and G.~Neven,
``Multi-Signatures in the Plain Public-Key Model and a General Forking Lemma,''
\textit{ACM CCS}, 2006.

\bibitem{rfc6979}
T.~Pornin,
``RFC 6979: Deterministic Usage of the Digital Signature Algorithm (DSA) and
Elliptic Curve Digital Signature Algorithm (ECDSA),''
IETF, 2013.
\url{https://www.rfc-editor.org/rfc/rfc6979}

\bibitem{wormhole2023}
Wormhole Foundation,
``Wormhole: A Generic Message Passing Protocol,''
2023.
\url{https://wormhole.com/whitepaper}

\bibitem{tbtc2022}
Keep Network,
``tBTC v2: A Decentralized Bitcoin Bridge,''
2022.
\url{https://docs.threshold.network/applications/tbtc-v2}

\bibitem{fireblocks2023}
Fireblocks,
``MPC-CMP: Threshold ECDSA from ECDSA Assumptions,''
2023.
\url{https://www.fireblocks.com/what-is-mpc/}

\bibitem{zengo2019}
ZenGo,
``Threshold Signatures: The Future of Private Keys,''
2019.
\url{https://zengo.com/mpc-wallet/}

\end{thebibliography}

